\subsection{BERNOULLI POLYNOMIALS}

Consider the composition operator U defined by

\[
U (f (\mathrm{X})) = f (\mathrm{X} + 1) - f (\mathrm{X});
\]

one can write \(U = e^{\mathrm{D}} - 1\) (VI, p.~270, Example 1); this is an operator of order 1, and if one puts \(U = DV\) one has \(V^{-1} = \frac{\mathrm{D}}{e^{\mathrm{D}} - 1}\) . The Appell polynomial of degree \(n\) corresponding to the operator \(U\) is called the Bernoulli polynomial of degree \(n\) and is denoted by \(\mathrm{B}_n(\mathrm{X})\) ; if one puts \(b_n = \mathrm{B}_n(0)\) one has the formulae

\[
\mathrm{B} _ {n} (\mathrm{X}) = \sum_ {k = 0} ^ {n} \binom{n}{k} b _ {n - k} \mathrm{X} ^ {k}\tag{18}
\]

\[
\frac {\mathrm{Se} ^ {\mathrm{XS}}}{e ^ {\mathrm{S}} - 1} = \sum_ {n = 0} ^ {\infty} \frac {1}{n !} \mathrm{B} _ {n} (\mathrm{X}) \mathrm{S} ^ {n}\tag{19}
\]

and in particular

\[
{\frac {\mathrm{S}}{e ^ {\mathrm{S}} - 1}} = \sum_ {n = 0} ^ {\infty} {\frac {1}{n !}} b _ {n} \mathrm{S} ^ {n}\tag{20}
\]

The formulae (7) and (9) of VI, p.~273, give, for the Bernoulli polynomials, the relations

\[
\frac {d \mathrm{B} _ {n}}{d \mathrm{X}} = n \mathrm{B} _ {n - 1} (X)\tag{21}
\]

\[
\mathrm{B} _ {n} (\mathrm{X} + 1) - \mathrm{B} _ {n} (\mathrm{X}) = n \mathrm{X} ^ {n - 1}.\tag{22}
\]

In particular, one has \(\mathrm{B}_{n}(1)-\mathrm{B}_{n}(0)=0\) for n\textgreater1, which, taking account of (18), gives the induction relation

\[
\sum_ {m = 0} ^ {n - 1} \binom {n} {m} b _ {m} = 0 \qquad (\text { for } n > 1)\tag{23}
\]

which enables one to calculate the \(b_{n}\) step-by-step. These numbers are clearly rational; since one can write

\[
{\frac {S}{e ^ {S} - 1}} = - {\frac {S}{2}} + {\frac {S}{2}} {\frac {e ^ {S} + 1}{e ^ {S} - 1}}
\]

and since (VI, p.~272, formula (5))

\[
\frac {e ^ {- S} + 1}{e ^ {- S} - 1} = - \frac {e ^ {S} + 1}{e ^ {S} - 1}
\]

one sees that in the formal series for \(\frac{S}{2}\frac{e^{S}+1}{e^{S}-1}\) all the terms of odd degree have coefficient zero; so one has

\[
b _ {0} = 1, \qquad b _ {1} = - \frac {1}{2}, \qquad b _ {2 n - 1} = 0 \quad \text { for } n > 1.\tag{24}
\]

The rational numbers \(b_{2n}\) ( \(n \geqslant 1\) ) are called the Bernoulli numbers; we shall see (VI, p.~288) that \(b_{2n}\) has the sign of \((-1)^{n-1}\) . The formula (23) gives, for the first values of n,

\[
\begin{array}{c} b _ {2} = \frac {1}{6}, \quad b _ {4} = - \frac {1}{3 0}, \quad b _ {6} = \frac {1}{4 2}, \quad b _ {8} = - \frac {1}{3 0}, \\ b _ {1 0} = \frac {5}{6 6}, \quad b _ {1 2} = - \frac {6 9 1}{2 7 3 0}, \quad b _ {1 4} = \frac {7}{6}, \\ b _ {1 6} = - \frac {3 6 1 7}{5 1 0}, \quad b _ {1 8} = \frac {4 3 8 6 7}{7 9 8}, \quad b _ {2 0} = - \frac {1 7 4 6 1 1}{3 3 0}, \quad b _ {2 2} = \frac {8 5 4 5 1 3}{1 3 8}, \\ b _ {2 4} = - \frac {2 3 6 3 6 4 0 9 1}{2 7 3 0}, \quad b _ {2 6} = \frac {8 5 5 3 1 0 3}{6}, \quad b _ {2 8} = - \frac {2 3 7 4 9 4 6 1 0 2 9}{8 7 0}. \end{array}
\]

Note that the numerators 691, 3617, 43867 are prime; the prime factorisations of the others are

\[
\begin{array}{c} 1 7 4 6 1 1 = 2 8 3 \times 6 1 7 \\ 8 5 4 5 1 3 = 1 1 \times 1 3 1 \times 5 9 3 \\ 2 3 6 3 6 4 0 9 1 = 1 0 3 \times 2 2 9 4 7 9 7 \\ 8 5 5 3 1 0 3 = 1 3 \times 6 5 7 9 3 1 \\ 2 3 7 4 9 4 6 1 0 2 9 = 7 \times 9 3 4 9 \times 3 6 2 9 0 3. \end{array}
\]

From these we deduce expressions for the first Bernoulli polynomials

\[
\mathrm{B} _ {0} (\mathrm{X}) = 1, \quad \mathrm{B} _ {1} (\mathrm{X}) = \mathrm{X} - \frac {1}{2}, \quad \mathrm{B} _ {2} (\mathrm{X}) = \mathrm{X} ^ {2} - \mathrm{X} + \frac {1}{6},
\]

\[
\mathrm{B} _ {3} (\mathrm{X}) = \mathrm{X} ^ {3} - \frac {3}{2} \mathrm{X} ^ {2} + \frac {1}{2} \mathrm{X}, \quad \mathrm{B} _ {4} (\mathrm{X}) = \mathrm{X} ^ {4} - 2 \mathrm{X} ^ {3} + \mathrm{X} ^ {2} - \frac {1}{3 0}.
\]

